\section{机械振动}
\sourceproblem{选择1}{12}{18}
\problemtext 比较弹跳、大角度单摆、矩形木块及球形木块的浮动，判断简谐运动。
\workings 矩形木块竖直浮动时，若保持部分浸没且水线截面积 $S$ 恒定，偏离平衡位置 $x$ 后合力为 $-\rho gSx$，故 $\ddot x=-(\rho gS/m)x$，选C。
\crosscheck 忽略阻力和流体附加质量；大角度单摆的回复力非线性，球形木块的水线面积通常变化。
\sourceproblem{选择2}{12}{18}
\problemtext 图1水平弹簧向左压缩后静止释放，右向为正，求余弦振动初相。
\workings $x(0)=-A,v(0)=0$，由 $x=A\cos(\omega t+\varphi)$ 得 $\varphi=\pi\pmod{2\pi}$，选D。
\crosscheck 初始回复力向右，初加速度为正。
\sourceproblem{选择3}{12}{18}
\problemtext 相同轻弹簧分别系不同质量及振幅的物体，比较水平、竖直振动周期。
\workings 理想线性弹簧 $T=2\pi\sqrt{m/k}$，重力只改变竖直振动的平衡位置，振幅不改变周期；仅B的质量不同，选B。
\crosscheck 质量倍增使周期乘 $\sqrt2$，并非乘2。
\sourceproblem{选择4}{12}{18}
\problemtext 从平衡点到半振幅位置的最短时间为 $t_0$，求周期。
\workings 令初速度朝向目标，$x=A\sin\omega t$，首次达到 $A/2$ 时 $\omega t_0=\pi/6$，故 $T=12t_0$，选B。
\crosscheck 最短时间小于四分之一周期。
\sourceproblem{选择5}{12}{18}
\problemtext 月面重力为地面的六分之一，同一小角度单摆周期如何变化？
\workings $T=2\pi\sqrt{l/g}$，故 $T_{\rm Moon}=\sqrt6T$，选C。
\crosscheck 摆长相同、小角近似成立；周期与质量无关。
\sourceproblem{选择6}{12}{18}
\problemtext 依据图2旋转矢量，初相 $\pi/4$、逆时针转过 $\pi t$，求振动方程。
\workings 取水平投影，$x=R\cos(\pi t+\pi/4)$，选A。
\crosscheck 初位置 $R/\sqrt2$、初速度 $-\pi R/\sqrt2$，与矢量转向一致。图在PDF19。
\sourceproblem{选择7}{12}{18}
\problemtext 依据图3，振幅10米，初位移负5米且上升，4秒处向下通过平衡点，求振动方程。
\workings 余弦初相取 $-2\pi/3$；图示首次向下过零对应相位 $\pi/2$，故 $4\omega=7\pi/6$，$\omega=7\pi/24$。$x=10\cos(7\pi t/24-2\pi/3)$（SI），选C。
\crosscheck 初速度为正，4秒处速度为负；根据PDF19图3排除多周期分支。
\sourceproblem{选择8}{13}{19}
\problemtext 图4振幅6米、周期4秒，2秒时向正方向过平衡点，求速度。
\workings $\omega=\pi/2\,\mathrm{rad/s}$，平衡点速率最大，$v(2)=A\omega=3\pi\,\mathrm{m/s}$，选A。
\crosscheck 原图从0时向负方向运动，2秒后正向过零。
\sourceproblem{选择9}{13}{19}
\problemtext 轻弹簧截成相等两段再并联，保持振动总能量不变，求新振幅。
\workings 均匀弹簧每半段劲度为 $2k$，并联为 $4k$。由 $kA^2/2=4kA'^2/2$ 得 $A'=A/2$，选A。
\crosscheck 弹簧均匀、质量忽略且满足线性关系；总能量相等是题设而非改装过程自动守恒。
\sourceproblem{选择10}{13}{19}
\problemtext 同向同频振动 $x_1=6\cos(5t+\pi/2)$、$x_2=2\cos(5t-\pi/2)$，求合振动。
\workings 两振动反相，$x=-6\sin5t+2\sin5t=-4\sin5t=4\cos(5t+\pi/2)$，选D。
\crosscheck 合振幅为差值4，频率不变。
\sourceproblem{选择11}{13}{19}
\problemtext 同频同向振动初相分别为 $\pi/3,-\pi/6$，振幅 $A_1,A_2$，求合振幅。
\workings 相差 $\pi/2$，矢量合成得 $A=\sqrt{A_1^2+A_2^2+2A_1A_2\cos(\pi/2)}=\sqrt{A_1^2+A_2^2}$，选C。
\crosscheck 任一分振幅为零时恢复另一分振幅。
